%==============================================================
\section{El Problema de los Datos: Una Perspectiva Histórica}
%==============================================================

Desde los albores de la civilización, la humanidad ha necesitado
almacenar y recuperar información de manera sistemática.
Las tablillas de arcilla sumerias (~3000 a.C.), los registros
contables del Imperio Romano, o los libros de contabilidad de los
comerciantes medievales son, en esencia, \textit{sistemas de gestión
de información}: colecciones de datos organizados con un propósito
definido y un mecanismo de consulta.

La revolución informática de mediados del siglo XX transformó este
problema milenario en una disciplina formal.
Con la proliferación de computadoras en organizaciones bancarias,
gubernamentales y empresariales durante los años 1950 y 1960,
la necesidad de almacenar grandes volúmenes de datos de forma
eficiente y recuperarlos con rapidez se volvió crítica.

\subsection{Era pre-relacional (1950–1970)}

Los primeros sistemas de gestión de datos surgieron como soluciones
ad-hoc a problemas específicos.
Predominaban dos modelos principales:

\begin{itemize}
  \item \textbf{Modelo jerárquico:} Los datos se organizaban en una
    estructura de árbol. IBM desarrolló IMS
    (\textit{Information Management System}) en 1966 para la misión
    Apolo de la NASA~\cite{date2003introduction}.
    Aunque eficiente para consultas predefinidas, era rígido:
    agregar una nueva relación entre datos requería reestructurar
    físicamente el sistema.

  \item \textbf{Modelo de red (CODASYL):} Propuesto por el comité
    CODASYL en 1969, permitía relaciones muchos-a-muchos mediante
    punteros físicos entre registros~\cite{codasyl1971}.
    Era más flexible que el jerárquico, pero extremadamente complejo
    de administrar y altamente dependiente de la estructura física
    del almacenamiento.
\end{itemize}

El problema fundamental de ambos modelos era la
\textbf{dependencia físico-lógica}: los programas de aplicación
debían conocer cómo estaban organizados físicamente los datos en
disco para poder acceder a ellos. Un cambio en la estructura de
almacenamiento implicaba reescribir todos los programas.

%==============================================================
\section{Edgar F.\ Codd y el Modelo Relacional}
%==============================================================

En junio de 1970, Edgar Frank Codd —matemático e investigador
de IBM— publicó el artículo que cambiaría para siempre el mundo de
las bases de datos:
\textit{``A Relational Model of Data for Large Shared Data Banks''}
en la revista \textit{Communications of the ACM}~\cite{codd1970}.

\begin{definicion}{Modelo Relacional Codd, 1970}
Un modelo de datos en el que la información se representa
exclusivamente mediante \textbf{relaciones} (tablas bidimensionales),
y en el que todas las operaciones sobre los datos se expresan
mediante un conjunto cerrado de operaciones del
\textbf{álgebra relacional}, sin necesidad de conocer la
organización física del almacenamiento.
\end{definicion}

La contribución de Codd fue revolucionaria por tres razones:

\begin{enumerate}
  \item \textbf{Independencia de datos:} Los programas de aplicación
    operan sobre una representación lógica (tablas), completamente
    desacoplada de cómo los datos se almacenan físicamente.

  \item \textbf{Fundamento matemático:} El modelo se basa en la
    teoría de conjuntos y la lógica de predicados de primer orden,
    lo que permite razonar formalmente sobre las operaciones.

  \item \textbf{Lenguaje declarativo:} El programador especifica
    \emph{qué} datos quiere, no \emph{cómo} obtenerlos.
    El sistema optimiza el acceso automáticamente.
\end{enumerate}

\subsection{Las 12 reglas de Codd}

En 1985, Codd formalizó su modelo en 12 reglas
(numeradas del 0 al 12) que debe cumplir un sistema para
ser considerado verdaderamente relacional~\cite{codd1985}.
Las más relevantes para este curso son:

\begin{table}[H]
\centering
\small
\renewcommand{\arraystretch}{1.3}
\begin{tabularx}{\columnwidth}{@{} c X @{}}
\toprule
\textbf{Regla} & \textbf{Descripción} \\
\midrule
0 & Gestión mediante capacidades relacionales \\
1 & Representación de la información como tablas \\
2 & Acceso garantizado a cada dato (tabla + clave + columna) \\
3 & Tratamiento sistemático de valores nulos \\
6 & Actualización de vistas \\
9 & Independencia de datos física y lógica \\
12 & No subversión del modelo relacional \\
\bottomrule
\end{tabularx}
\caption{Selección de las 12 reglas de Codd (1985).}
\label{tab:reglas-codd}
\end{table}

%==============================================================
\section{Nacimiento de SQL: De SEQUEL a SQL}
%==============================================================

A raíz de la publicación de Codd, IBM inició el proyecto
\textit{System R} en 1973 en sus laboratorios de San José,
California, con el objetivo de demostrar la viabilidad práctica del
modelo relacional~\cite{chamberlin1974}.
Donald D.\ Chamberlin y Raymond F.\ Boyce diseñaron un lenguaje de
consulta denominado \textbf{SEQUEL}
(\textit{Structured English Query Language}).

El nombre fue posteriormente acortado a \textbf{SQL}
(\textit{Structured Query Language}) por razones legales de marca
registrada.
En paralelo, Michael Stonebraker y Eugene Wong desarrollaron
\textbf{INGRES} en la Universidad de California, Berkeley (1973–1977),
que introdujo su propio lenguaje de consulta llamado QUEL.

\begin{nota}
Aunque coloquialmente se pronuncia tanto ``sequel'' como
``ese-cu-ele'', el estándar ISO/IEC lo denomina formalmente
``SQL''. En este curso usaremos la pronunciación ``ese-cu-ele''
por su precisión técnica.
\end{nota}

\subsection{El primer sistema comercial: Oracle}

En 1979, la empresa Relational Software Inc.\ (hoy Oracle
Corporation), fundada por Larry Ellison, liberó el primer
sistema de gestión de bases de datos relacionales comercial:
\textbf{Oracle V2}.
Paradójicamente, Oracle llegó al mercado antes que el propio
IBM con su DB2 (1983), a pesar de que la tecnología subyacente
provenía de los laboratorios de IBM.


%==============================================================
\section{¿Qué es una Base de Datos?}
%==============================================================

\begin{definicion}{Base de Datos}
Una \textbf{base de datos} es una colección organizada, persistente
y compartible de datos relacionados entre sí, diseñada para
minimizar la redundancia y maximizar la integridad, de modo que
múltiples aplicaciones y usuarios puedan acceder a ella de manera
concurrente y controlada~\cite{date2003introduction}.
\end{definicion}

Esta definición merece descomponerse en sus atributos esenciales:

\begin{itemize}
  \item \textbf{Organizada:} Los datos siguen una estructura formal
    (esquema) que define qué tipos de información se almacenan y
    cómo se relacionan entre sí.

  \item \textbf{Persistente:} Los datos sobreviven a la terminación
    de los procesos que los crean. A diferencia de las variables en
    memoria, los datos en una base de datos permanecen aunque el
    sistema se apague.

  \item \textbf{Compartible:} Múltiples usuarios o aplicaciones
    pueden acceder a los mismos datos simultáneamente, de forma
    controlada y sin interferencias.

  \item \textbf{Mínima redundancia:} El mismo dato no se repite
    innecesariamente en múltiples lugares, lo que reduce el riesgo
    de inconsistencias (un dato cambia en un lugar pero no en otro).
\end{itemize}

\subsection{Analogía con el mundo real}

Imaginemos el sistema de una biblioteca universitaria:

\begin{ejemplo}
Una biblioteca almacena información sobre \textbf{libros},
\textbf{estudiantes} y \textbf{préstamos}.
Sin una base de datos, cada ficha de préstamo copiaría el título
completo del libro, el autor, el nombre del estudiante, su carrera,
etc.
Si un estudiante cambia su dirección, habría que actualizar decenas
de fichas físicas.

Con una base de datos relacional, el estudiante aparece
\emph{una sola vez} en la tabla \texttt{Estudiantes}.
Todos los préstamos simplemente hacen referencia a su identificador.
Un cambio de dirección se realiza en un único lugar y se refleja
automáticamente en todo el sistema.
\end{ejemplo}




